\begin{abstract}
High--resolution rotational spectroscopy provides extremely precise
information on molecular structure and rovibrational dynamics.
The parameters entering Watson's effective rotational Hamiltonian,
such as quartic and sextic centrifugal distortion constants, have
traditionally been determined from spectroscopic fits of observed
transition frequencies, whereas modern quantum-chemical calculations
can access the same quantities through direct evaluation of
rovibrational couplings.  Establishing a transparent connection
between these inverse and direct viewpoints has therefore become
increasingly important for the interpretation of modern spectroscopic
data.

In this work we develop a tensor-based framework for the analysis and
transformation of centrifugal distortion parameters.  The tensorial
structure itself is already implicit in classical rovibrational
perturbation theory; the present contribution is to make it explicit
as the basis of a practical framework that treats spectroscopic fits
and direct quantum-mechanical calculations on the same footing.
For the quartic case, the natural reduced object is a symmetric
$3\times3$ tensor $\Tau'$, from which the conventional Watson
constants are obtained by a unique forward projection.  Conversely,
the reconstruction of the tensor from a set of Watson constants is an
affine gauge problem, for which the Moore--Penrose pseudoinverse
selects a distinguished representative.

Based on this viewpoint, we introduce a computational procedure that
reconstructs reduced quartic and sextic tensor representatives
directly from spectroscopic parameter sets.  Representation
transformations are then obtained by simple permutations of tensor
indices followed by reprojection of the transformed tensor onto the
desired Watson reduction and principal-axis representation.  For
quartic centrifugal distortion the tensor algorithm reproduces the
analytical transformations of Yamada and Winnewisser with machine
precision, while for sextic centrifugal distortion it provides a
practical and internally consistent extension for cases where
analytical representation transforms are not generally available.

The resulting formulation provides a unified and numerically stable
framework for handling centrifugal distortion parameters obtained from
spectroscopic fits, quantum-chemical calculations, or rovibrational
computations, and offers a representation-independent tensorial
perspective on centrifugal distortion in molecular spectroscopy.
\end{abstract}
